\cvsection{Open-Source and Artifacts}
\cvitem{\textbf{Certified language model robustness:} \href{https://github.com/UCF-ML-Research/CR-UTP}{CR-UTP}.}
\cvitem{\textbf{Backdoor detection and mitigation:} \href{https://github.com/UCF-ML-Research/SSL-Cleanse}{SSL-Cleanse}.}
\cvitem{\textbf{Externally checked programming agents:} \href{https://fhe-coder.github.io/}{FHE-Coder benchmark and project resources}.}
\cvitem{\textbf{Trustworthy learning and private computing:} \href{https://github.com/UCF-ML-Research}{UCF-ML-Research}; \href{https://github.com/UCF-Lou-Lab-PET}{UCF-Lou-Lab-PET}.}
\cvsection{Selected Research Contributions and Impact}
\begin{itemize}
\item \textbf{Attacks and red teaming.} TrojLLM (NeurIPS 2023) and TrojText (ICLR 2023) expose Trojan risks in language models. MemIncept (ICML 2026) and conjunctive prompt attacks (ACL 2026) examine agent memory and interaction. RealM (NeurIPS 2026) provides a red-teaming benchmark for physical-world vision-language models.
\item \textbf{Defenses and reliable agents.} CR-UTP (Findings of ACL 2024) certifies robustness within a specified text-perturbation setting. SSL-Cleanse (ECCV 2024) detects and mitigates Trojans in self-supervised learning. FHE-Coder (ICLR 2026) uses functional and security checks to guide cryptographic programming agents.
\item \textbf{Provenance and evaluation.} PRO (NeurIPS 2026) studies robust text watermarking; SIF (CVPR 2026) develops fingerprints for vision-language models. EMNLP 2025/2026 studies evaluate watermarking under cross-lingual changes and in medical text, examining factuality and usefulness alongside detectability.
\item \textbf{Private and verifiable computation.} SecureRoute (DAC 2026) protects model selection and inference. DataSeal (IEEE S\&P 2025) checks supported encrypted computations. These mechanisms address confidentiality and computation integrity, complementing rather than replacing behavioral safety measures.
\end{itemize}
